\documentclass[11pt,reqno]{amsart}

\usepackage[margin=1in]{geometry}
\usepackage{amsmath,amssymb,amsthm,mathrsfs}
\usepackage{enumitem}
\usepackage{xcolor}
\usepackage{hyperref}
\usepackage[capitalise,noabbrev]{cleveref}

\definecolor{darklink}{rgb}{0.0,0.2,0.5}
\hypersetup{
  colorlinks=true,
  linkcolor=darklink,
  citecolor=darklink,
  urlcolor=darklink
}

\theoremstyle{plain}
\newtheorem{theorem}{Theorem}[section]
\newtheorem{lemma}[theorem]{Lemma}
\newtheorem{proposition}[theorem]{Proposition}
\newtheorem{corollary}[theorem]{Corollary}
\newtheorem{conjecture}[theorem]{Conjecture}

\theoremstyle{definition}
\newtheorem{definition}[theorem]{Definition}
\newtheorem{example}[theorem]{Example}
\newtheorem{problem}[theorem]{Open problem}

\theoremstyle{remark}
\newtheorem{remark}[theorem]{Remark}

\newcommand{\Channel}{\mathbf{Channel}}
\newcommand{\Sigmaspace}{\mathbf{\Sigma}}
\newcommand{\PCF}{\mathrm{PCF}}
\newcommand{\Lt}{L(t)}
\newcommand{\BoT}{\mathrm{BoT}}
\newcommand{\CC}{\mathrm{CC}}
\newcommand{\Vquad}{V_{\mathrm{quad}}}
\newcommand{\dn}{\delta_{n}}
\newcommand{\Borel}{\mathcal{B}}
\newcommand{\Pade}{\mathrm{Pad\acute e}}

\title[Channel structure of asymptotic deformations in PCF dynamics]
{Channel structure of asymptotic deformations\\ in polynomial continued fraction dynamics:\\ an outline}

\author{The SIARC author}
\date{2026-05-01 (draft)}

\begin{document}
\maketitle

\begin{abstract}
We propose, define, and catalogue the notion of an
\emph{asymptotic channel} for sequences arising from polynomial
continued fractions (PCFs). A channel is a triple $(D,T,S)$
consisting of a dense subspace $D$ of formal series, an
asymptotic-gauge topology $T$ on $D$, and a section $S$ from
analytic germs to $D$. Three concrete channels appear in the
SIARC stack: the recurrence-parameter channel $\Lt$, the
Borel-of-trans-series channel $\BoT$, and the
connection-coefficient channel $\CC$. We summarise the empirical
record (six $\Delta<0$ degree-$2$ families surveyed in three
recent computational sessions), state a conjectural no-go
theorem for $\Lt$ and a speculative bridge identity between
$\Lt$ and $\BoT$, and identify the questions that the SIARC
Master Conjecture's channel functor
$\chi : \Sigmaspace \to \Channel$ must answer before the
announcement paper can be written rigorously. This is an outline
and program statement: no new theorems are proved.
\end{abstract}

\tableofcontents

% =====================================================================
\section{Position}
\label{sec:position}

The SIARC stack has accumulated three empirically distinct
phenomena that all advertise themselves as Painlev\'e
reductions:
\begin{enumerate}[label=(\alph*)]
\item The \emph{Stokes signal} $S(t)<1$ measured in the
recurrence-parameter channel for the six known $\Delta<0$
degree-$2$ polynomial continued fractions
(\cite{siarc_pcf1_v13}, Sec.~5).
\item The explicit reduction
$\Vquad \rightsquigarrow P_{III}(D_6)$ at the parameter point
$(1/6, 0, 0, -1/2)$, established in the connection-coefficient
channel via the V$_\mathrm{quad}$ Casoratian / Heun bridge
(\cite{siarc_pcf1_v13}, Sec.~6;
\cite{slavyanov2000special, iwasaki1991from}).
\item The total failure, demonstrated in two computational
sessions of 2026-05-01 (\cite{siarc_painleve_deep_6x,
siarc_borel_5x, siarc_borel_kscan}), of both the $\Lt$ channel
and a 1/$n$ trans-series Borel-resummation channel to recover
the $\Vquad$ reduction. Painlev\'e ansatz residuals stay flat
under precision escalation (no-fit), and Pad\'e--Borel pole
radii drift unboundedly under truncation extension.
\end{enumerate}
Item~(c) is critical. It rules out the most naive interpretation
of the Master Conjecture's \emph{channel functor}
$\chi : \Sigmaspace \to \Channel$ \cite{siarc_umbrella}, namely
that any deformation channel detects any reduction. The Master
Conjecture asserts that the channel is a function of the
deformation $\Sigma$, not of the family. The empirical record
of (a)--(c) shows that $\chi$ must in fact be non-trivial:
$\chi(\Sigma_{\Lt}) \ne \chi(\Sigma_{\CC})$, and the two
channels carry genuinely different analytic data about the
same underlying sequence.

This outline catalogues what we know about the channel space,
states the questions $\chi$ must answer, and isolates a small
number of conjectural identities whose resolution would
constitute the technical content of the announcement paper.

\paragraph{Scope.}
We restrict throughout to PCF families
\[
  \dn = b_n - \cfrac{a_n}{b_{n-1} - \cfrac{a_{n-1}}{\ddots - a_1/b_0}}
\]
in the Lorentzen--Waadeland sense \cite{lorentzen2008continued},
with $a_n, b_n$ polynomial in $n$ over $\mathbf{Z}$, $b$
$\mathbf{Z}$-primitive and irreducible, and (in the empirical
sections) $\deg b = 2$ with discriminant $\Delta < 0$.

\paragraph{What is and is not in scope of this outline.}
This is a program statement and a definition document. We
(i)~fix terminology, (ii)~enumerate the channels that the
SIARC stack has actually built, (iii)~report the empirical
verdict on each, and (iv)~state the open problems that the
Master Conjecture's channel functor must resolve. We do
\emph{not} prove a no-go theorem (\Cref{conj:nogo} is
conjectural), we do \emph{not} construct the bridge identity
(\Cref{conj:bridge} is speculative), and we do \emph{not}
extend the catalogue to $\deg b \ge 3$ or to families with
$\Delta \geq 0$. Each of those is a self-contained follow-up
listed in \Cref{sec:open}.

\paragraph{Reading guide.}
\Cref{sec:def}--\Cref{sec:catalogue} are definitional and
should be read by anyone who wants to know what a SIARC
``channel'' is. \Cref{sec:functor} packages the Master
Conjecture in channel language. \Cref{sec:nogo}--%
\Cref{sec:bridge} are the two non-trivial conjectures whose
resolution is the core technical content of the Master
Conjecture announcement. \Cref{sec:implications}--%
\Cref{sec:open} convert the conjectures into a sequenced
research program.

% =====================================================================
\section{Definition: asymptotic channel}
\label{sec:def}

\begin{definition}[Channel]\label{def:channel}
Let $X = (\dn)_{n \geq n_0}$ be a sequence of real or
complex numbers of subexponential growth. An
\emph{asymptotic channel} for $X$ is a triple $(D, T, S)$
where
\begin{enumerate}
\item $D$ is a dense subspace of a fixed completion
$\widehat{D}$ of formal power series (or trans-series)
in a distinguished variable $z$, equipped with a fixed
trans-series scaffold (the \emph{template});
\item $T$ is a Hausdorff topology on $D$ induced by an
asymptotic gauge $\Gamma : D \to \mathbf{R}_{\ge 0}^{\mathbf{N}}$
that compares partial sums against a chosen scale;
\item $S : \mathcal{O}_0 \dashrightarrow D$ is a partially
defined section of the natural inclusion of analytic
germs at $z=0$ into $\widehat{D}$. We write
$\mathrm{dom}\, S \subset \mathcal{O}_0$.
\end{enumerate}
The sequence $X$ is said to \emph{lie in the channel}
$(D,T,S)$ if there exists a germ $f \in \mathrm{dom}\, S$
such that $\Gamma(\widetilde X - S(f))$ is summable in the
sense of $T$, where $\widetilde X \in \widehat D$ is a
canonically associated formal expansion of $X$.
\end{definition}

The triple $(D,T,S)$ encodes three independent choices:
\emph{which asymptotic ansatz to fit} ($D$),
\emph{how to measure goodness of fit} ($T$), and
\emph{which analytic continuation rule to invert} ($S$).
Different choices yield different channels even for the same
underlying sequence.

\begin{example}[Recurrence-parameter channel $\Lt$]
\label{ex:lt}
$D$ is the space of formal series in $1/n$ with coefficients
analytic in a deformation parameter $t \in (-h_0, h_0)$, with
the Birkhoff--Trjitzinsky template
$\dn(t) \sim \exp(\rho \log n) \cdot n^\sigma \cdot
\sum_{k \ge 0} c_k(t)/n^k$
\cite{birkhofftrjitzinsky1932analytic, wimp1984computation}.
The gauge $\Gamma$ is uniform convergence on partial sums of
length $N$, evaluated on a finite $t$-grid of step $h$. The
section $S$ is Borel--Laplace summation along $\mathbf{R}_+$
\cite{costin2008asymptotics}. This is the channel of
\cite{siarc_pcf1_v13}, Sec.~5.
\end{example}

\begin{example}[Borel-of-trans-series channel $\BoT$]
\label{ex:bot}
$D$ is the space of trans-series
$\log|\dn| \sim -A n \log n + \alpha n - \beta \log n
+ \gamma + \sum_{k \ge 1} h_k / n^k$,
with the WKB scaffold $(A,\alpha,\beta,\gamma)$ fitted first
and then frozen. The gauge $T$ measures Pad\'e convergence
of the Borel transform
$\Borel(\zeta) = \sum_{k \ge 1} h_k \zeta^{k-1}/(k-1)!$,
quantified by stability of the principal pole radius
$|\rho_0|$ under truncation order $K$. The section $S$ is
Pad\'e--Borel resummation. This is the channel of
\cite{siarc_borel_5x, siarc_borel_kscan}.
\end{example}

\begin{example}[Connection-coefficient channel $\CC$]
\label{ex:cc}
Suppose $\dn$ is the second solution of a second-order linear
difference (or differential) equation with regular singular
points. $D$ is the space of formal series in a transfer-matrix
variable $\tau$ around a chosen singular point. The gauge
$T$ is convergence of Stokes constants extracted via
isomonodromic deformation \cite{jimbomiwa1981monodromy}. The
section $S$ is the Riemann--Hilbert correspondence:
analytic germs at the singular point are reconstructed from
the connection matrix
\cite{lodayrichaud2016divergent, boalch2001stokes}. The
established $\Vquad \rightsquigarrow P_{III}(D_6)$ at
$(1/6,0,0,-1/2)$ \cite{iwasaki1991from, ohyama2006studies}
lives in this channel.
\end{example}

\begin{remark}[Channel vs.\ summation method]
\label{rem:channel-vs-sum}
At first sight a channel is just a choice of summation method
(Borel--Laplace vs.\ Pad\'e--Borel vs.\ isomonodromic
reconstruction). The distinction we draw is finer: two
channels can use the same summation method (e.g.\ both $\Lt$
and $\BoT$ ultimately appeal to a Borel sum) yet target
different formal-series spaces $D$ and use different
asymptotic gauges $T$. The empirical content of
\Cref{sec:position}(c) is precisely that this distinction
is not vacuous on the data side: $\Lt$ and $\BoT$ both fail
where $\CC$ succeeds, despite all three being candidate
inverses of "what $\dn$ is asymptotically doing".
\end{remark}

\begin{remark}[Why this is not a tautology]
\label{rem:not-tautology}
The definition admits the trivial channel $(D,T,S)$ with $D$
full, $T$ discrete, and $S$ the identity on its domain, in
which every sequence ``lies''. The non-trivial content of
the channel framework is that the natural channels arising
from (i)~asymptotic enumeration of recurrences, (ii)~Borel
resummation of trans-series, and (iii)~isomonodromic Stokes
analysis are mutually \emph{non-cofinal}: the partial order
on channels by ``inverse-image inclusion'' (i.e.\
$(D,T,S) \le (D',T',S')$ iff every sequence in the first
channel lies in the second) does not flatten the three
working channels of \Cref{ssec:lt}--\Cref{ssec:cc} into a
single one. This non-flattening is the empirical observation
that motivates the bridge conjecture (\Cref{conj:bridge})
and its potential failure mode.
\end{remark}

\begin{remark}[Trans-series template as a choice]
\label{rem:template}
The Birkhoff--Trjitzinsky template
$\dn(t) \sim e^{\rho \log n} n^{\sigma} \sum_k c_k(t)/n^k$
is distinguished but not unique. One could equally well
take a $\sqrt n$-template (which Session E's iteration v2
tried, before being abandoned;
\cite{siarc_borel_5x}) or a hybrid template
mixing exponential and Stokes factors. Each yields a
different channel by \Cref{def:channel}. The empirical
failure of the $\sqrt n$-template for $\Vquad$ shows that
channel choice is not free: a wrong template gives a
well-defined but vacuous channel.
\end{remark}

% =====================================================================
\section{Catalogue of channels in the SIARC stack}
\label{sec:catalogue}

We summarise the three channels in turn. Numerical values are
the published values; sources are cited in line.

\subsection{The recurrence-parameter channel $\Lt$}
\label{ssec:lt}

\paragraph{Definition.} \Cref{ex:lt}.

\paragraph{Operationalisation.} For each PCF family, perturb
$b_n \mapsto b_n + t \cdot \omega_n$ with a chosen direction
$\omega_n$ (typically a constant-term or root-radius shift),
compute $\dn(t)$ on a finite $t$-grid via exact symbolic
recursion, and fit a Birkhoff--Trjitzinsky tail.

\paragraph{Empirical findings.}
\begin{itemize}
\item \emph{Stokes signal.} For all six known $\Delta<0$
degree-$2$ families
$\{Q_{L01}, Q_{L02}, Q_{L06}, \Vquad, Q_{L15}, Q_{L26}\}$,
the Stokes proxy $S(t)$ enters the strip $S<1$ after
sufficient depth, with stability across precision sweeps
\cite{siarc_pcf1_v13}.
\item \emph{Painlev\'e ansatz fit (Session D).} 0/6 families
admit a Painlev\'e reduction at residual $\le 10^{-4}$
under depth $3000$, dps $400$. $\Vquad$ marks marginal at
$4.59 \times 10^{-5}$ on a $P_{III}$ template but does
\emph{not} recover the known $(1/6, 0, 0, -1/2)$ parameters
\cite{siarc_painleve_deep_6x}.
\item \emph{Channel-typing observation.} $Q_{L06}$ and
$Q_{L26}$ both lie in $\mathbf{Q}(\sqrt{-7})$ but disagree
on the best Painlev\'e ansatz ($P_V$ vs.\ $P_{III}$), so
CM-field-driven Painlev\'e typing is ruled out at the $\Lt$
level.
\end{itemize}

\paragraph{Open implementation questions.}
\begin{itemize}
\item Does enriching the deformation direction $\omega_n$ from
constant-term to a moduli of perturbations (e.g.\ a smooth
path in coefficient space) move any family into a positive
$\Lt$ Painlev\'e fit?
\item Is the $S<1$ signal a Painlev\'e fingerprint at all,
or is it (as Sessions D suggest) an artefact of a coarser
analytic structure that does not reach the
connection-coefficient resolution?
\end{itemize}

\subsection{The Borel-of-trans-series channel $\BoT$}
\label{ssec:bot}

\paragraph{Definition.} \Cref{ex:bot}.

\paragraph{Operationalisation.} Strip the $-A n \log n$
factor with the closed-form WKB exponents
\[
  A = \begin{cases} 4 & a_n \text{ constant}, \\
                    3 & a_n = c_a n + d_a, \end{cases}
  \qquad
  \alpha \;=\; A - 2 \log c_b + \log|c_a|
\]
verified to $\ge 13$ digits across all six families
\cite{siarc_borel_kscan}. Fit the residual trans-series
$h_k/n^k$ by least squares, then compute Pad\'e approximants
$[L/M]$ to the Borel transform $\Borel(\zeta)$ and probe the
nearest singularity.

\paragraph{Empirical findings.}
\begin{itemize}
\item \emph{WKB exponent identity.} 6/6 families to
$\ge 13$ digits; this is locked into PCF-1 v1.3 and is the
only positive structural result of Sessions D--E$'$.
\item \emph{Borel-resummation gate ($\Vquad$).} Painlev\'e
fit residuals on the Pad\'e--Borel sum
\emph{worsen} under truncation extension
($K=12 \to 24$): $7.2\times 10^{-3} \to 5.0 \times 10^{-3}
\to 1.0 \times 10^{-2} \to 9.9 \times 10^{-2}$. Pad\'e pole
radius diverges $2.17 \to 4.59$. Best-fit Painlev\'e identity
flips $P_{III} \leftrightarrow P_{VI}$
\cite{siarc_borel_5x}.
\item \emph{$K$-scan on marginal hits.} Two $K=12$ marginal
hits (Q$_{L15} \to P_V$ at $9 \times 10^{-5}$,
Q$_{L26} \to P_{III}$ at $3 \times 10^{-5}$) likewise drift
under $K$-extension: residuals fail to stay $\le 10^{-4}$,
Pad\'e pole radii drift unboundedly, ansatz identity flips
\cite{siarc_borel_kscan}. Confirmed Pad\'e-truncation
artefacts.
\end{itemize}

\paragraph{Open implementation questions.}
\begin{itemize}
\item Does a different summation scheme (Borel--\'Ecalle
\cite{ecalle1981fonctions, sauzin2014resurgent} rather than
Pad\'e--Borel) on the \emph{same} trans-series space $D$
recover Painlev\'e information for $\Vquad$? If yes, the
"channel" we tested was actually $(D,T,S')$ with the wrong
section $S'$, not the wrong space $D$.
\item Are the $h_k$ stable enough across $K$ to attempt a
Borel--Laplace sum along a non-real ray? (Sessions E--E$'$
report $h_1 \ldots h_6$ stable to 14 digits across
$K=12$--$24$; the channel definition is well-posed even
though the downstream sum fails.)
\end{itemize}

\subsection{The connection-coefficient channel $\CC$}
\label{ssec:cc}

\paragraph{Definition.} \Cref{ex:cc}.

\paragraph{Operationalisation.} For families that arise as
specialisations of a Heun-type linear ODE
\cite{slavyanov2000special}, compute the isomonodromic
deformation parameters and read off the Painlev\'e moduli
point.

\paragraph{Empirical findings.}
\begin{itemize}
\item $\Vquad$: $P_{III}(D_6)$ at $(1/6, 0, 0, -1/2)$
\cite{iwasaki1991from, ohyama2006studies}.
\item Other five $\Delta<0$ degree-$2$ families: pipeline
not yet executed (op:cc-pipeline below).
\end{itemize}

\paragraph{Open implementation questions.}
\begin{itemize}
\item Build an explicit $\CC$-channel pipeline for the five
non-$\Vquad$ families. The pipeline requires identifying
the underlying second-order linear difference equation,
computing its Stokes data, and matching against the
$P_{II}$--$P_{VI}$ catalogue.
\item Does the $\CC$ pipeline succeed on \emph{all} of the
five remaining families, on a strict subset, or on none?
Each outcome implies a different shape of the channel
functor $\chi$.
\end{itemize}

\subsection{Comparison table}
\label{ssec:comparison}

The three channels can be summarised in a single table.
``Detects $\Vquad$'s $P_{III}(D_6)$ reduction'' is the
natural test question; the answer is the empirical content
of Sessions D, E, E$'$.

\begin{center}
\renewcommand{\arraystretch}{1.25}
\begin{tabular}{lcccc}
\hline
Channel & Space $D$ & Gauge $T$ & Section $S$ &
  Detects $\Vquad$? \\
\hline
$\Lt$  & series in $1/n$, $t$-deformed &
  uniform on $t$-grid &
  Borel--Laplace & no \\
$\BoT$ & 1/$n$ trans-series of $\log|\dn|$ &
  Pad\'e pole-radius stability &
  Pad\'e--Borel & no \\
$\CC$  & series in $\tau$ at sing.\ point &
  Stokes-constant convergence &
  Riemann--Hilbert & yes \\
\hline
\end{tabular}
\end{center}

\noindent The empirical asymmetry in the rightmost column is
the single-line summary of the channel-functor problem: a
rigorous $\chi$ must explain why exactly the $\CC$ row
detects this reduction and the others do not.

\subsection{The locked WKB exponent identity}
\label{ssec:wkb}

A byproduct of the $\BoT$ channel construction
(\Cref{ssec:bot}) is the explicit closed form for the WKB
scaffold $(A, \alpha)$ extracted from $\log |\dn|$. We
record the value because it is the only positive structural
output of Sessions D--E$'$ and because it appears in PCF-1
v1.3 \cite{siarc_pcf1_v13}.

\begin{proposition}[WKB exponent identity, empirical]
\label{prop:wkb}
For each of the six $\Delta<0$ degree-$2$ families surveyed,
the leading trans-series exponents of
$\log|\dn| = -A n \log n + \alpha n - \beta \log n + \gamma
+ \sum_{k \ge 1} h_k / n^k$
satisfy the closed form
\[
  A \;=\; \begin{cases} 4 & a_n \equiv \text{const}, \\
                        3 & a_n = c_a n + d_a, \end{cases}
  \qquad
  \alpha \;=\; A \;-\; 2 \log c_b \;+\; \log|c_a|
\]
to $\geq 13$ decimal digits across the $K$-scan grid
$K \in \{12, 16, 20, 24\}$, dps $\in \{2200, 3000, 4000,
5000\}$ \cite{siarc_borel_kscan}.
\end{proposition}

\begin{center}
\renewcommand{\arraystretch}{1.2}
\begin{tabular}{lccccc}
\hline
Family & $\Delta$ & $a_n$ & $A$ & $\alpha$ &
  digits matched \\
\hline
Q$_{L01}$ & $-3$  & const & $4$ & $4$              & exact \\
Q$_{L02}$ & $-4$  & const & $4$ & $4 + \log 2$     & 13 \\
Q$_{L06}$ & $-7$  & lin   & $3$ & $3 - 2\log 2$    & 14 \\
$\Vquad$  & $-11$ & const & $4$ & $4 - 2\log 3$    & 13 \\
Q$_{L15}$ & $-20$ & const & $4$ & $4 - 2\log 3$    & 15 \\
Q$_{L26}$ & $-28$ & const & $4$ & $4 + \log 3 - 2\log 4$ & 13 \\
\hline
\end{tabular}
\end{center}

\Cref{prop:wkb} is what survives of the $\BoT$ channel after
the Pad\'e--Borel resummation gate fails: the channel's
\emph{template} is well-conditioned, the $h_k$ are stable to
$14$ digits across $K=12$--$24$ \cite{siarc_borel_kscan},
but the downstream resummation does not converge to a
well-defined section $S$. In the language of
\Cref{def:channel}: $D$ and $T$ are well-defined; $S$ is
not. This is the most refined diagnostic available for
what goes wrong in $\BoT$.

% =====================================================================
\section{The channel functor: what we know and what we don't}
\label{sec:functor}

The Master Conjecture \cite{siarc_umbrella} posits a functor
\[
  \chi : \Sigmaspace \longrightarrow \Channel,
\]
where $\Sigmaspace$ is a category of deformations of PCF
families (objects: families with chosen deformation direction;
morphisms: compatible reparametrisations) and $\Channel$ is
the category whose objects are channels in the sense of
\Cref{def:channel}.

\subsection{The deformation category $\Sigmaspace$}
\label{ssec:sigma}

We make $\Sigmaspace$ slightly more precise, since the
channel functor is only as well-defined as its source.

\begin{definition}[Object of $\Sigmaspace$]
An object of $\Sigmaspace$ is a triple
$\Sigma = (F, \omega, h)$ where $F = \PCF(a, b)$ is a PCF
family as in the scope of \Cref{sec:position}, $\omega \in
\mathbf{Z}[n]$ is a deformation direction, and $h > 0$ is a
step parameter. The deformed family is
$F_t = \PCF(a, b + t \omega)$ for $|t| < h$.
\end{definition}

\begin{definition}[Morphism of $\Sigmaspace$]
A morphism
$(F, \omega, h) \to (F', \omega', h')$ exists when $F = F'$,
$\omega' = \lambda \omega$ for $\lambda \in \mathbf{Q}^\times$,
and $h' \le h / |\lambda|$. Composition is multiplication of
the scalars $\lambda$. The resulting category is a Picard
groupoid over each fixed family.
\end{definition}

The channel functor $\chi$ is required to factor through the
family: $\chi(F, \omega, h)$ is independent of $h$ and
invariant under positive rescalings of $\omega$. This makes
$\chi$ a functor on the quotient $\Sigmaspace / \sim$, but
we retain the unquotiented form because the choice of
$\omega$ class \emph{does} affect $\chi$ (different
perturbation directions can in principle produce different
channels).

\subsection{Properties of $\chi$}

\begin{conjecture}[Properties of $\chi$]
\label{conj:chi}
The channel functor satisfies:
\begin{enumerate}[label=(\roman*)]
\item \emph{Functoriality in $\Sigmaspace$.} Changing the
deformation $\Sigma$ on a fixed family genuinely changes the
target channel $\chi(\Sigma)$. Equivalently, $\chi$ is not
constant in its first argument.
\item \emph{Empirical consistency.} For every PCF family
$F$ for which a Painlev\'e reduction is presently established
(at this writing: $\Vquad$ alone, in the $\CC$ channel),
$\chi$ assigns the deformation that produced that reduction
to the channel in which it was established. Concretely:
$\chi(\Sigma(\Vquad)) = \CC$.
\item \emph{Non-constancy.} There exist deformations
$\Sigma_1, \Sigma_2$ with $\chi(\Sigma_1) \ne \chi(\Sigma_2)$.
\end{enumerate}
\end{conjecture}

We comment on each clause.

\medskip
(i) is testable: any positive $\Lt$ Painlev\'e fit in any
new family would furnish an instance of $\chi$ taking value
$\Lt$, distinguishing it from $\chi(\Sigma(\Vquad)) = \CC$.

\medskip
(ii) is currently a single data point. The cross-validation
record (\Cref{sec:position}(c)) shows that for $\Vquad$ the
$\Lt$ and $\BoT$ channels do not return the $\CC$ reduction;
this is consistent with (ii) but does not yet test it on
multiple families.

\medskip
(iii) is asserted by (a)--(c) of \Cref{sec:position} together
with $\Vquad$'s $\CC$ reduction.

\paragraph{What we do not know about $\Channel$.}
The most basic open question is whether the channel space is
discrete or has moduli structure. The three channels
$\{\Lt, \BoT, \CC\}$ exhausted in this outline are
clearly not the only ones in principle: any choice of
$(D, T, S)$ defines a channel. Whether the \emph{useful} set
of channels (those that detect a non-trivial structure on
some PCF family) is finite, countable, or has continuous
moduli is open, and is the content of the
\textsf{(op:channel-moduli)} problem in \Cref{sec:open}.

\paragraph{Worked instance of $\chi$ on the present catalogue.}
The minimal information required from $\chi$ on the data
collected so far is the assignment
\[
  \chi : \big( \Vquad,\ \text{Heun-isomonodromic}\big)
  \;\longmapsto\; \CC,
  \qquad
  \chi : \big( F_i,\ \Sigma_{\text{const-term}}\big)
  \;\longmapsto\; ?,
\]
for $i = 1, \ldots, 6$ and the $\Lt$-style constant-term
deformation $\Sigma_{\text{const-term}}$. Sessions D and
E establish that the second assignment is \emph{not} $\Lt$
(no Painlev\'e ansatz reaches residual $\le 10^{-4}$ stably)
and \emph{not} $\BoT$ (Pad\'e pole radii drift). Either
$\chi$ takes value in some channel not yet in our catalogue,
or the deformation $\Sigma_{\text{const-term}}$ is
``channel-degenerate'' in the sense that no Painlev\'e
structure is detectable along it for these families.
\Cref{conj:nogo} below conjectures the latter for $\deg
b = 2$, $\Delta < 0$; \Cref{conj:bridge} indicates how
the two readings could be reconciled via a bridge identity.

% =====================================================================
\section{Conjectured no-go theorem}
\label{sec:nogo}

\begin{conjecture}[$\Lt$ no-go for degree $2$, $\Delta<0$]
\label{conj:nogo}
Let $F = \PCF(a, b)$ be a polynomial continued fraction with
$a, b \in \mathbf{Z}[n]$, $b$ irreducible and
$\mathbf{Z}$-primitive, $\deg b = 2$, and discriminant
$\Delta(b) < 0$. Then no Painlev\'e reduction of the
$\Lt$-deformation
$F_t = \PCF(a, b + t\omega)$ exists in the
$\Lt$ channel of \Cref{ex:lt}.
\end{conjecture}

\paragraph{Empirical evidence.} \Cref{conj:nogo} is consistent
with the Session D record \cite{siarc_painleve_deep_6x}:
0 / 6 families admit a Painlev\'e fit at residual
$\le 10^{-4}$, with residuals stable under precision
escalation (the diagnostic that the candidate ODE is genuinely
wrong rather than under-resolved).

\paragraph{Falsifiability.} \Cref{conj:nogo} is falsified by a
single counterexample: any $\Delta<0$ degree-$2$ family for
which an $\Lt$-channel Painlev\'e fit residual drops below
$10^{-4}$ under precision escalation, with stable
parameters.

\paragraph{Failure modes to monitor.}
\begin{itemize}
\item Non-generic perturbation directions $\omega$. Sessions
B--D used D-A (constant-term) and D-B (root-radius). A
generic moduli of $\omega$ has not been swept.
\item Higher-rank deformations (two-parameter $t$). Out of
scope here; would change the formal-series space $D$.
\item Families with reducible Galois group of $b$. Out of
scope by the irreducibility hypothesis, but worth flagging
as a sanity boundary (cf.\ Ap\'ery in
\cite{siarc_pcf2_program}).
\end{itemize}

% =====================================================================
\section{Conjectured bridge identity}
\label{sec:bridge}

\begin{conjecture}[Bridge identity, speculative]
\label{conj:bridge}
There exists a real-analytic function
$\Phi(s, \zeta_1, \ldots, \zeta_r)$ such that, for every
$\Delta<0$ degree-$2$ PCF family $F$ in the regime of
\Cref{conj:nogo}:
\begin{enumerate}[label=(\roman*)]
\item $s$ is the limiting Stokes exponent
$S(t) = \lim_{t \to 0^+} S_t$ measured in the $\Lt$
channel,
\item $\zeta_1, \ldots, \zeta_r$ are the singular points of
the Borel transform $\Borel(\zeta)$ in the $\BoT$ channel
(if these are well-defined; presently they are not),
\item $\Phi(s, \zeta_1, \ldots, \zeta_r) = 0$ holds, and
\item $\Phi$-membership of $F$ is equivalent to existence of
a $\CC$-channel Painlev\'e reduction of $F$.
\end{enumerate}
\end{conjecture}

\Cref{conj:bridge} is the speculative deep claim. If true, the
three channels $\Lt$, $\BoT$, $\CC$ are not analytically
independent: they parameterise the same underlying structure
along three different probes, and the Stokes signal $S(t)<1$
in $\Lt$ is a coarse shadow of a sharper rigidity that lives
in $\CC$.

\paragraph{Status.} We have one data point in the affirmative
direction: $\Vquad$ does have $S < 1$ in $\Lt$
\cite{siarc_pcf1_v13} and \emph{does} have a $\CC$-channel
$P_{III}(D_6)$ reduction \cite{iwasaki1991from}. The other
five families have $S<1$ but no known $\CC$ reduction yet.
\Cref{conj:bridge} predicts they all have one.

\paragraph{Obstructions.} The main obstruction to even
\emph{stating} \Cref{conj:bridge} precisely is that
$\zeta_1, \ldots, \zeta_r$ are not well-defined for the five
non-$\Vquad$ families: Sessions E and E$'$ exhibited unbounded
Pad\'e pole drift \cite{siarc_borel_5x, siarc_borel_kscan}.
Either the trans-series template is wrong, the Pad\'e
approximants are inadequate (Borel--\'Ecalle along a complex
ray would be the next test), or the $\BoT$ channel is genuinely
empty for these families and the bridge identity is between
$\Lt$ and $\CC$ alone, with $\BoT$ playing no role.

The connection-formula viewpoint of
\cite{costincostin2001formation} gives the most natural
formal home for $\Phi$: $\Phi$ would be a relation among
Stokes constants attached to two different singular-point
neighbourhoods of the same difference equation. Whether such
a relation can be written down explicitly for any $\Delta<0$
PCF is the technical content of (op:bridge-vquad).

\paragraph{Three structural variants of \Cref{conj:bridge}.}
We distinguish three weakened forms, in increasing order of
falsifiability:
\begin{enumerate}[label=(B\arabic*)]
\item \emph{Existence.} Some real-analytic $\Phi$ as in
\Cref{conj:bridge} exists. (The default reading.)
\item \emph{Algebraicity.} $\Phi$ is a polynomial in $s$ with
coefficients algebraic in $\zeta_1, \ldots, \zeta_r$. This
is the strongest reasonable form and would constrain the
channel functor to take values in a discrete set of
algebraic strata.
\item \emph{Approximate.} For some prescribed tolerance
$\varepsilon$, $|\Phi(s, \vec\zeta)| \le \varepsilon$ on a
dense subset of $\Delta < 0$ degree-$2$ families. This is
the weakest form that still gives empirical traction: it
would reduce $\chi$ from a deterministic functor to an
$\varepsilon$-coarse functor and would still permit a
statistical announcement.
\end{enumerate}
The rubber-duck critique flagged that without committing
to one of (B1)--(B3), \Cref{conj:bridge} is too elastic to
be falsifiable. We commit to (B1) as the \emph{minimum}
requirement and to (B2) as the \emph{aspiration}; (B3) is the
fallback for the announcement paper if (B1)/(B2) prove
infeasible within the relay timeframe.

\paragraph{Sketch of $\Phi$ for $\Vquad$.}
For $\Vquad$ the $\CC$-channel reduction parameters
$(1/6, 0, 0, -1/2)$ in $P_{III}(D_6)$ are exactly determined
(via the Heun pathway of \cite{slavyanov2000special} and the
IKSY classification \cite{iwasaki1991from, ohyama2006studies}).
From \Cref{prop:wkb}, $A = 4$ and
$\alpha = 4 - 2 \log 3$. The $\Lt$ Stokes proxy $S(t)$
measured in \cite{siarc_pcf1_v13} satisfies $S < 1$
unconditionally. A candidate skeleton for $\Phi$ would
therefore relate $S$, $A$, $\alpha$, and a (currently
undefined) Borel singular point $\zeta_0$ via
\[
  \Phi_{V} (s, \zeta_0) \;\stackrel{?}{=}\;
  s + A^{-1}\bigl(\alpha - 2\log\,c_b\bigr)
  - \Re(\zeta_0)/A.
\]
We stress that this is a structural placeholder, not a
claim: the right-hand side is gauge-fixed to vanish for
$\Vquad$ by construction, and its empirical content lives
in whether the same shape with $A, \alpha$ replaced by the
row values of \Cref{prop:wkb} produces a $\zeta_0$ that
actually appears as a Borel singularity for the other five
families. As of Session E$'$, no such $\zeta_0$ exists at
the Pad\'e level, so the sketch is currently untestable.

% =====================================================================
\section{Related work and prior channel structures}
\label{sec:related}

The channel framework of \Cref{def:channel} did not arise in
a vacuum. We sketch four lines of prior work that supply
either the formal vocabulary or a non-PCF analogue, and we
flag where the SIARC channels diverge from each.

\subsection{Resurgence and alien calculus
(\'Ecalle, Costin, Sauzin)}
The most developed pre-existing analogue is the
resurgence-theoretic framework of
\cite{ecalle1981fonctions, costin2008asymptotics,
sauzin2014resurgent, mitschi2016divergent}. There a
divergent series is identified with its Borel transform,
analytic continuation along admissible rays defines a
collection of summed functions, and the differences between
those sums (Stokes constants) carry the deep arithmetic
content. Our $\BoT$ channel is essentially a discrete
Pad\'e-truncated version of this construction; the failure
mode observed in Session E
(\cite{siarc_borel_5x}) is consistent with the well-known
phenomenon that Pad\'e approximation does not always recover
the alien structure even when the underlying series is
resurgent. The continuous-side success of resurgence
\cite{delabaere2016divergent} on $P_I$ is the structural
benchmark to which any positive resolution of
\textsf{(op:bridge-vquad)} must aspire.

\subsection{Stokes / isomonodromic theory
(Birkhoff, Jimbo--Miwa, Boalch)}
The connection-coefficient channel $\CC$ has a long history
in the theory of isomonodromic deformations of linear ODEs
\cite{jimbomiwa1981monodromy, vanderputsinger2003galois}.
Boalch's work \cite{boalch2001stokes} packages the Stokes
data as a Poisson--Lie structure on a moduli space, which is
the natural ambient candidate for the channel-moduli question
\textsf{(op:channel-moduli)}. The non-classical aspect for
PCFs is that the underlying object is a second-order linear
\emph{difference} (not differential) equation; the
isomonodromic theory transfers but is not entirely
streamlined.

\subsection{Heun--Painlev\'e correspondence
(Slavyanov--Lay)}
\cite{slavyanov2000special} treats Heun-class linear ODEs
with four regular singular points and their reductions to
the Painlev\'e family. The $\Vquad$ Casoratian arises as a
specialisation of this framework, and that specialisation is
how the $(1/6, 0, 0, -1/2)$ moduli point is computed in
practice. Whether the other five $\Delta < 0$ degree-$2$ PCFs
admit Heun-style specialisations of any kind is an empirical
question that \textsf{(op:cc-pipeline)} would answer.

\subsection{Algebraic Painlev\'e moduli
(Mazzocco, IKSY, Ohyama et al.)}
\cite{mazzocco2001picard, mazzocco2002painleve, iwasaki1991from,
ohyama2006studies, guzzetticlass2012} give finite catalogues of
algebraic and Picard-type solutions to the Painlev\'e
equations together with the Stokes-data parameters that
realise them. Our T2B / UMB-PVI-MATCH umbrella result
\cite{siarc_t2b_v3} is an attempt to read PCF families
\emph{into} this catalogue: Class B (16/16) lands inside
$P_{III}(D_6)$, and a single $P_{VI}$-Picard rigidity
candidate (A052) was flagged for further screening. That
candidate, if it survives, will be the second example
beyond $\Vquad$ where $\chi$ takes value in $\CC$ on a
specific PCF.

\medskip
The novel content of the present outline is not in the
formal-series machinery, which is well established, but in
the \emph{empirical} observation that on the natural
$\Delta < 0$ degree-$2$ PCF stratum the three classical
channels are \emph{empirically separated} on the
$\Vquad$-detection task. Whether this separation is
intrinsic to PCFs or an artefact of the channels we have
chosen to test is the central open question.

% =====================================================================
\section{Implications for the Master Conjecture}
\label{sec:implications}

A rigorous announcement of the Master Conjecture
\cite{siarc_umbrella} requires the following intermediate
results, in increasing order of difficulty:

\begin{enumerate}
\item \textbf{Pipeline closure.} An explicit $\CC$-channel
pipeline for the five non-$\Vquad$ degree-$2$ $\Delta<0$
families (\textsf{op:cc-pipeline}). This is the
\emph{minimum} test of $\chi$: \Cref{conj:chi}(ii) demands
that whatever the pipeline returns is the value of $\chi$ on
those families' deformations.
\item \textbf{No-go proof or refutation.} \Cref{conj:nogo}
either proved (settling that $\Lt$ never detects Painlev\'e
in this regime) or refuted by a positive counterexample
(\textsf{op:no-go-proof}). Either outcome sharpens
\Cref{conj:chi}(i).
\item \textbf{Partial bridge identity.} A version of
\Cref{conj:bridge} for at least one family beyond $\Vquad$
(\textsf{op:bridge-vquad}). This is the first non-trivial
test of the speculative bridge.
\item \textbf{Cubic extension.} Channel catalogue extended to
$\deg b \ge 3$ along the lines of \cite{siarc_pcf2_program}
(\textsf{op:cubic-channels}).
\end{enumerate}

In the language of \Cref{sec:functor}: items 1 and 2
discharge \Cref{conj:chi}(i)--(ii) for one stratum
($\deg = 2$, $\Delta < 0$); item 3 starts to characterise the
target category $\Channel$; item 4 begins to extend $\chi$
across the natural product structure of $\Sigmaspace$.

\medskip
The announcement paper cannot be written rigorously without
\emph{at least} item 1 above (the $\CC$ pipeline), because
otherwise the empirical content of \Cref{conj:chi}(ii)
collapses to a single example.

% =====================================================================
\section{Open problems and program}
\label{sec:open}

We collect six concrete open problems with tractability
estimates. "Session" here means one SIARC relay session
(of order $10$ hours of agent compute plus human-side review).

\begin{problem}[\textsf{op:no-go-proof}]
Prove or refute \Cref{conj:nogo}.
\emph{Tractability.} A refutation may be within reach in
1--3 sessions if any moduli of perturbation directions
$\omega$ produces a stable $\Lt$ Painlev\'e fit. A proof is a
multi-year program: it requires structural understanding of
why the Birkhoff--Trjitzinsky template's gauge cannot
distinguish Painlev\'e equations within the $\Delta<0$
regime.
\end{problem}

\begin{problem}[\textsf{op:bridge-vquad}]
Construct a candidate $\Phi$ in \Cref{conj:bridge} for
$\Vquad$ alone, using its known $\CC$ reduction as the
target.
\emph{Tractability.} 2--5 sessions for a candidate, if a
suitable Borel--\'Ecalle resummation can be implemented
along a complex ray; multi-year for a proof.
\end{problem}

\begin{problem}[\textsf{op:cc-pipeline}]
Implement an explicit $\CC$-channel pipeline for the five
non-$\Vquad$ degree-$2$ $\Delta<0$ families. Report which (if
any) admit a Painlev\'e reduction.
\emph{Tractability.} 1--2 sessions per family for the
straightforward Heun-type cases (cf.\ V\_quad's
\cite{slavyanov2000special} pathway); harder if the
underlying linear difference equation does not arise from
classical Heun specialisation.
\end{problem}

\begin{problem}[\textsf{op:channel-moduli}]
Is the channel space $\Channel$ a moduli space (with
positive-dimensional families of channels) or a discrete /
finite set? Equivalently, does varying $(D,T,S)$ continuously
through admissible triples give continuous deformations of
the channel?
\emph{Tractability.} Multi-year. Likely requires the
Stokes-matrix / Frobenius-manifold framework of
\cite{boalch2001stokes} as a guide, and a working
$\CC$ pipeline as an empirical anchor.
\end{problem}

\begin{problem}[\textsf{op:cubic-channels}]
Do PCF-2 cubic families (\cite{siarc_pcf2_program}) admit the
same three-channel structure $\{\Lt, \BoT, \CC\}$, or does
$\deg b = 3$ introduce qualitatively new channels?
\emph{Tractability.} 3--5 sessions for a first scan once the
PCF-2 enumeration cubic Apery-positive-control pipeline is
in place.
\end{problem}

\begin{problem}[\textsf{op:partition-link}, speculative]
Does the Meinardus partition-asymptotic universality of G-01
\cite{meinardus1954asymptotische, andrews1976theory} live in
any of the channels of \Cref{sec:catalogue}? If so, the
channel functor $\chi$ extends across SIARC Path 5 and unifies
the partition and PCF strands.
\emph{Tractability.} Multi-year. Requires reformulating
Meinardus universality as a channel-detected property of a
suitable trans-series.
\end{problem}

\begin{problem}[\textsf{op:de-borel}]
Implement Borel--\'Ecalle ray summation
\cite{ecalle1981fonctions, sauzin2014resurgent} as a
\emph{distinct} section $S'$ on the same trans-series
template $D$ as $\BoT$. Test on $\Vquad$. If $S'$ recovers
the $\CC$ moduli point, the SIARC catalogue gains a fourth
channel $\BoT' = (D, T', S')$ and the empirical asymmetry of
\Cref{ssec:comparison} sharpens to a statement about Pad\'e
gauges specifically, not about Borel resummation in
general. \emph{Tractability.} 2--4 sessions; well within
relay scope, and likely the highest-value next experiment.
\end{problem}

\begin{problem}[\textsf{op:zero-one-law}]
Is there a zero-one law for channel detection on a fixed
PCF family? Concretely: given a family $F$, is the set of
deformations $\Sigma$ for which $\chi(\Sigma) \ne \mathsf{trivial}$
either of measure zero or of full measure in a natural
sense? A positive answer would force $\chi$ to factor through
a discrete invariant of the family.
\emph{Tractability.} Multi-year and currently lacks a
candidate measure on $\Sigmaspace$.
\end{problem}

% =====================================================================
\section*{Acknowledgements}

We thank the SIARC relay infrastructure
(\href{https://github.com/papanokechi/siarc-relay-bridge}%
{papanokechi/siarc-relay-bridge}) for storing and
checkpointing the empirical record cited throughout.

\paragraph{AI disclosure.} This outline was drafted with
the assistance of large language models in the SIARC
self-iterating analytic relay chain. All numerical claims
trace to exact computations stored under
\texttt{sessions/2026-05-01/} on the SIARC bridge. The
present document contains no new theorems and no new
numerical results; it organises and frames existing
empirical and structural material into a program statement.

% =====================================================================
\bibliographystyle{plain}
\bibliography{annotated_bibliography}

\end{document}
